% !TeX program = xelatex
% 主入口。UTF-8 编码。构建：latexmk -xelatex main.tex
% 基于 caowenhan/hhuthesis v2.0.0 的明确标识适配类。
% 暂按“电子信息”预设 professionalmaster；本人核实后可改 academicmaster
% 或 nonfulltimemaster。不要把本稿的暂设视为已确认学位/培养方式。
\IfFileExists{hhuthesis.cls}{%
  \documentclass[professionalmaster]{hhuthesis}% 与完整原仓库合并时自动使用原类
}{%
  \documentclass[professionalmaster]{hhuthesis-cad}% 本交付包的独立编译入口
}
% ===== 写作阶段开关 =====
\newif\ifoutlineguidance
\outlineguidancetrue  % 改为 \outlineguidancefalse 可隐藏“写作提示”。
\newif\ifdraftlayout
\draftlayouttrue      % 写作预览不插空白页；定稿前改为 false 并核对学院要求。
\newif\iffullfrontmatter
\fullfrontmatterfalse % 当前仅摘要/大纲，默认不输出正式授权声明和全套封面。

% 三级目录：章 -> 节 -> 小节，例如 4 / 4.3 / 4.3.2。
\setcounter{secnumdepth}{2}
\setcounter{tocdepth}{2}
\titlecontents{subsection}[2\ccwd]{\songti\zihao{-4}}
  {\thecontentslabel\hspace{0.5\ccwd}}{}
  {\hspace{.5em}\titlerule*{.}\contentspage}
\hypersetup{hidelinks,pdfdisplaydoctitle=true,
 pdftitle={异构CAD形变特征转换技术研究与模型转换实现：摘要与三级目录初稿},
 pdfauthor={作者待填写},pdfsubject={硕士论文写作框架}}
\emergencystretch=2em
\raggedbottom
% 非正式预览采用连续起章，保留模板的其余页边距、字体级别及页眉页脚。
\ifdraftlayout
 \makeatletter\@openrightfalse\makeatother
 \let\cleardoublepage\clearpage
 % 仅提纲预览压缩标题留白；关闭 draftlayout 后恢复原模板间距。
 \ctexset{section={beforeskip=10pt,afterskip=8pt},subsection={beforeskip=8pt,afterskip=6pt}}
\fi
% 这些提示是提纲说明，不是最终论文正文。
\newcommand{\chapterguide}[1]{%
 \ifoutlineguidance\par\noindent{\small\color{black!70}\textbf{本章写作重点：}#1}\par\medskip\fi}
\newcommand{\outlinehint}[1]{%
 \ifoutlineguidance\par\noindent{\small\color{black!70}\textbf{写作提示：}#1}\par\fi}

% 所有“待填写”均需本人确认。未推断姓名、学号、导师、学院、答辩日期。
% 专业“电子信息”和题目来自中期汇报；学位类别在 main.tex 中暂设专业硕士。
\studentnumber{待填写}
\classification{待填写}
\securitylevel{待确认}
\udc{待填写}
\title{异构CAD形变特征转换技术研究\\与模型转换实现}
\vtitle{异构CAD形变特征转换技术研究与模型转换实现}
\author{待填写}
\tutorinfoa{待填写}
\tutorinfob{待填写}
\degree{学位名称待确认}
\major{电子信息}
\submitdate{待填写}
\defenddate{待填写}
\awarded{河海大学，日期待填写}
\chairman{待填写}
\reviewera{待填写}
\reviewerb{待填写}
\reviewerc{待填写}
\reviewerd{待填写}
\reviewere{待填写}
\reviewerf{}
\reviewerg{}
\nlcdate{待填写}
\nlclocate{中~~国~~$\cdot$~~南~~京}
\institute{河海大学}
\zhtitle{异构CAD形变特征转换技术研究与模型转换实现}
\zhsubtitle{无}
\entitle{Research on Parametric Form Feature Translation and Model Conversion across Heterogeneous CAD Systems}
\ensubtitle{None}
\thesislang{汉语}
\abstractlang{汉、英}
\thesispages{待填}
\numofwords{待填}
\thesiskeywords{异构CAD；统一特征模型；可编辑模型}
\researchfield{CAD特征转换（待本人确认）}
\tutor{待填写}
\tutorinstitute{待填写}
\tutora{待填写}
\tutorainstitute{待填写}
\tutorb{待填写}
\tutorbinstitute{待填写}
\englishtitle{Research on Parametric Form Feature Translation and Model Conversion across Heterogeneous CAD Systems}
\englishauthor{To be completed}
\entutor{To be completed}
\eninstitute{Hohai University}
\englishdepartment{Department to be completed}
\englishdate{Date to be completed}
\enlocate{Nanjing, P.R. China}
\englishdegree{Master's degree (to be confirmed)}

\begin{document}
\hypersetup{pageanchor=false}
\iffullfrontmatter
  % 正式封面接口保留；下列完整模式须在信息/格式核定后重新检查。
  \makenlctitle{}
  % \makeverticaltitle
  \makeinfo{}
  \makeeninfo{}
  \makedeclare
\else
  % 写作阶段说明封面，不冒充学校正式送审封面。
\begin{titlepage}
\thispagestyle{empty}
\begin{center}
\vspace*{16mm}
{\heiti\zihao{1}河海大学\par}
\vspace{8mm}
{\songti\zihao{2}硕士学位论文写作框架\par}
\vspace{22mm}
{\heiti\zihao{2}异构CAD形变特征转换技术研究\par 与模型转换实现\par}
\vspace{10mm}
{\songti\zihao{3}中英文摘要与三级目录初稿\par}
\vspace{18mm}
{\songti\zihao{4}
\begin{tabular}{rl}
研究范围：& 草图、基准面、拉伸、旋转、扫掠\\[6pt]
研究系统：& SolidWorks、Creo、NX\\[6pt]
项目基础：& CADExchange、CADProj\\[6pt]
专业名称：& 电子信息\\[6pt]
作者及导师：& 待填写
\end{tabular}\par}
\vfill
{\small 本稿包含拟定目录和写作提示，不是完成的学位论文。\par
实验数据、正式文献、本人信息及最终格式均需补齐和核对。\par}
\vspace{8mm}
{\songti\zihao{4}2026年10月24日\par}
\vspace*{8mm}
\end{center}
\end{titlepage}

\fi
\frontmatter
\hypersetup{pageanchor=true}
% 摘要初稿：无未经验证的实验数据。正文完成后补充经验证的结果并同步中英文。
\begin{abstract}
异构CAD系统之间的零件模型交换不仅需要保持几何形状，还需要保留特征参数和构造依赖，以支持模型在目标系统中的继续修改。针对SolidWorks、Creo和NX在参数含义、空间方向、对象组织及引用方式等方面的差异，本文围绕基于草图的形变特征开展转换研究，以拉伸、旋转和扫掠为核心，以草图与基准面为支撑，研究统一表达、参数映射和目标模型重建方法。主要研究内容如下。

首先，从几何构造与系统参数调查出发，归纳数值量纲、方向约定、终止条件、薄壁材料位置、草图约束和几何引用等差异，明确结构、参数、几何及编辑能力的保留要求。针对同名参数含义不同和相同外形对应不同构造关系的问题，以参数的实际作用作为统一表达的依据。

其次，设计由基础几何、单位、对象引用、草图与基准面、成形特征组成的统一模型。分别组织草图图元与约束作用位置、基准面构造条件、拉伸与旋转的两端范围及薄壁偏置，并用独立草图、内嵌草图和参数圆形截面表达扫掠的不同截面来源。通过标识和引用描述特征依赖，以必要几何信息辅助目标对象定位。

然后，研究面向目标原生特征的转换方法。通过源与目标实际草图坐标的转换恢复轮廓和约束方向，结合构造参考重建基准面；通过方向、终止范围和两侧材料位置的映射处理拉伸与旋转；通过截面来源、路径起点、曲线顺序和截面放置的联合处理实现基本扫掠重建，并明确辅助构造与受限参数的适用条件。

最后，基于CADExchange与CADProj实现公共模型层及三系统读写适配，形成模型读取、参数转换、顺序重建和结果记录流程。设计三条同系统基线和六条跨系统路线，从特征结构、参数与约束、几何一致性、目标原生编辑和执行效率等方面组织评价。现有工作形成了五类特征的统一表达与转换原型，复杂约束、特殊终止及高级扫掠选项仍需完善，相关保留程度应由明确范围内的实验分别验证。

\keywords{异构CAD；形变特征；统一特征模型；参数映射；引用恢复；可编辑模型}
\end{abstract}

\begin{enabstract}
Part-model exchange across heterogeneous CAD systems requires both geometric consistency and the preservation of parameters and dependencies for further editing. This thesis studies sketch-based parametric form feature translation across SolidWorks, Creo, and NX. Extrusion, revolution, and sweep are the principal features, supported by sketches and datum planes.

First, differences in units, directions, termination conditions, thin-wall material placement, sketch constraints, and references are analyzed through geometric construction and parameter investigation. Requirements for structure, parameters, geometry, and editability are distinguished. Parameters are interpreted by their actual construction effects rather than their names.

Second, a unified model organizes basic geometry, units, references, sketches, datum planes, and form features. It represents sketch entities and constraint attachment points, datum-plane construction conditions, two-sided extents, and thin-wall offsets. Sweep profiles are classified as independent sketches, embedded sketches, or parametric circular sections. Identifiers and references describe dependencies, while geometric information assists target-object location.

Third, reconstruction methods recover native features. Transformations between actual sketch frames restore profile geometry and constraint directions, while construction references guide datum-plane reconstruction. Direction, termination-range, and material-placement mappings support extrusion and revolution. Basic sweep reconstruction jointly handles profile sources, path starting points, segment order, and section placement, with explicit conditions for auxiliary constructions and restricted parameters.

Finally, CADExchange and CADProj implement the common model and system adapters for extraction, parameter conversion, sequential reconstruction, and result recording. An evaluation scheme uses three same-system baselines and six directed cross-system routes, separately examining feature structure, parameters and constraints, geometry, native editing, and efficiency. The work provides a unified representation and a conversion prototype for the five feature categories. Complex constraints, special termination cases, and advanced sweep options remain limited; preservation levels require experimental verification within explicit scopes.

\enkeywords{heterogeneous CAD; parametric form features; unified feature model; parameter mapping; reference recovery; editable models}
\end{enabstract}

\tableofcontents
% 正文有实际图表后再启用，避免在提纲阶段产生空清单。
% \listoffigures
% \listoftables
\mainmatter
\chapter{绪论}
\label{chap:1}
\chapterguide{说明研究动机、已有方法、研究范围和论文主线。研究现状需用正式阅读的文献补齐，不能把项目设计直接写成已证明的学术创新。}

\section{研究背景与意义}
\label{sec:1-1}

\subsection{异构CAD环境下的模型交换需求}
\label{subsec:1-1-1}
\outlinehint{从协同设计、软件迁移和历史零件复用引出跨系统交换需求，结合实际工作场景说明目标模型仍需继续修改。}

\subsection{参数化模型转换中的信息保留问题}
\label{subsec:1-1-2}
\outlinehint{区分外形、特征类型、尺寸参数、引用关系和修改后的更新行为；用定深切除与到面切除的区别说明仅外形一致仍不充分。}

\subsection{可编辑特征转换的研究意义}
\label{subsec:1-1-3}
\outlinehint{围绕减少重复建模和支持模型延续使用阐述意义。零件库检索、批量变体与训练数据只作为潜在用途，不作为已验证成果。}

\section{国内外研究现状}
\label{sec:1-2}

\subsection{几何与拓扑信息交换方法}
\label{subsec:1-2-1}
\outlinehint{调研中性格式和几何内核交换的实际信息范围，区分本项目使用的终态几何导出与标准本身可能支持的更广信息。}

\subsection{基于特征回放与宏参数化的转换方法}
\label{subsec:1-2-2}
\outlinehint{比较特征序列、参数和约束的提取与重建机制，关注映射范围、源目标依赖和参考对象处理方式。}

\subsection{基于统一模型的特征转换方法}
\label{subsec:1-2-3}
\outlinehint{比较中间表达与语义模型如何组织参数和依赖，讨论系统扩展、参数差异与目标重建的研究问题。}

\subsection{转换结果验证方法与现有研究不足}
\label{subsec:1-2-4}
\outlinehint{分别梳理结构、参数、几何和编辑验证方法，在真实文献对照后归纳不足，不把几何成功当作可编辑性已完整保留。}

\section{研究目标与主要研究内容}
\label{sec:1-3}

\subsection{形变特征定义与研究范围}
\label{subsec:1-3-1}
\outlinehint{沿用中期题目中的“形变特征”，明确指基于草图与参数驱动的成形特征，不指自由曲面或网格变形。以拉伸、旋转、扫掠为核心，以草图、基准面为支撑；轴、点、边、面仅作为相关引用涉及。}

\subsection{研究目标与关键问题}
\label{subsec:1-3-2}
\outlinehint{提出四个问题：参数实际含义如何统一，空间与引用如何恢复，公共表达如何转为目标建模步骤，以及保留程度如何评价。}

\subsection{主要研究任务与个人工作边界}
\label{subsec:1-3-3}
\outlinehint{组织差异分析、统一模型、转换算法、系统实现与实验任务；按实际分工说明扫掠的主导与协作部分，其他工程特征不列为个人核心贡献。}

\section{技术路线与论文结构}
\label{sec:1-4}

\subsection{总体技术路线}
\label{subsec:1-4-1}
\outlinehint{给出“几何构造分析—三系统参数调研—公共表达—目标重建—分层验证”的路线，并标明实验反馈如何修正字段和映射。}

\subsection{论文章节安排}
\label{subsec:1-4-2}
\outlinehint{说明第二章回答差异和需求，第三章定义保存什么，第四章解释怎样转换，第五章说明系统落地，第六章评价结果，第七章收敛结论。}

\chapter{异构CAD特征差异分析与转换需求}
\label{chap:2}
\chapterguide{按参数、空间、组织和依赖等维度归纳问题，并逐类落实到五类特征。此章分析“差异是什么”，不提前展开完整算法。}

\section{研究对象与差异分析框架}
\label{sec:2-1}

\subsection{三系统参数调查的对象与方法}
\label{subsec:2-1-1}
\outlinehint{以SolidWorks、Creo、NX的特征定义、读取参数和创建条件为调查对象，记录系统版本、数据来源和可访问范围。}

\subsection{九类异构问题的分类与对应关系}
\label{subsec:2-1-2}
\outlinehint{保留九类分析框架：量纲、方向、参数含义、对象组织、约束依赖、引用历史、几何拓扑、功能范围、执行环境；建立分类与五类特征的交叉表。}

\section{草图与基准面的表达差异}
\label{sec:2-2}

\subsection{草图坐标、图元与圆弧表达差异}
\label{subsec:2-2-1}
\outlinehint{比较局部坐标与世界坐标、圆弧范围、构造线以及椭圆读取和创建范围；通过最小案例说明直接复制数值的风险。}

\subsection{草图约束、尺寸与外部关联差异}
\label{subsec:2-2-2}
\outlinehint{比较约束参与对象、起终点和圆心等作用位置，区分驱动尺寸、几何关系和对外部几何的关联。}

\subsection{基准面构造方法与参考条件差异}
\label{subsec:2-2-3}
\outlinehint{比较偏置、角度、中间面、过点过线等构造和参考组合，说明相同平面位置不等于相同构造依赖。}

\section{拉伸、旋转与扫掠的参数差异}
\label{sec:2-3}

\subsection{拉伸方向、深度与终止条件差异}
\label{subsec:2-3-1}
\outlinehint{对比单向、双向、对称、到对象与到延伸面等含义，使用总深度60 mm与单侧30 mm说明换算需求。}

\subsection{旋转轴、转动方向与角度范围差异}
\label{subsec:2-3-2}
\outlinehint{区分轴线本身的方向和特征的正转方向，比较两侧角度的正负约定以及参考终止可能产生的互补角。}

\subsection{扫掠截面、路径与方向控制差异}
\label{subsec:2-3-3}
\outlinehint{比较独立草图、内嵌截面、参数圆截面，分析路径起点、曲线段顺序、截面放置和沿路径方向控制的共同影响。}

\section{跨特征的共性差异}
\label{sec:2-4}

\subsection{数值量纲、模板单位与薄壁材料位置差异}
\label{subsec:2-4-1}
\outlinehint{区分长度、角度与无量纲参数；比较厚度加方向与两侧偏置，说明总厚度相同仍可能形成不同的材料位置。}

\subsection{引用身份、建模历史与拓扑划分差异}
\label{subsec:2-4-2}
\outlinehint{说明原生编号不可跨系统复用，分析父特征归属变化、边分段与接缝对引用恢复和几何比较的影响。}

\subsection{目标功能范围与构造组织差异}
\label{subsec:2-4-3}
\outlinehint{归纳独立草图与内嵌截面、辅助基准面、子类型不对称、参数未使用和原生执行条件的差异。}

\section{转换需求与评价目标归纳}
\label{sec:2-5}

\subsection{公共表达与目标重建需求}
\label{subsec:2-5-1}
\outlinehint{从差异推导统一量纲、明确方向、带角色的引用、约束作用位置及可解释组合重建等需求，形成可追踪的需求表。}

\subsection{能力边界与结果评价需求}
\label{subsec:2-5-2}
\outlinehint{区分结构可表达、源端能读取、目标能创建和实验证明四个层次；提出结构、参数、几何和编辑再生的分别检查要求。}

\section{本章小结}
\label{sec:2-6}

\outlinehint{归纳本章能够支持的认识，并说明与下一章的衔接；不重复列举所有接口或扩展未证明结论。}

\chapter{面向五类特征的统一模型设计}
\label{chap:3}
\chapterguide{围绕参数及其作用设计中间结构，保持与现有CADExchange定义一致。可选扩展与尚未由三系统完整读写的字段要单独标注。}

\section{统一模型的设计原则与总体结构}
\label{sec:3-1}

\subsection{基于几何构造与参数调查的设计过程}
\label{subsec:3-1-1}
\outlinehint{按构造分析、系统参数调查、公共含义归纳和读写反馈设计结构，避免把原生接口字段机械取并集。}

\subsection{分层结构与特征历史组织}
\label{subsec:3-1-2}
\outlinehint{划分基础几何、标识引用、草图基准、成形特征和配套证据层；说明有序特征列表与引用的关系，不扩展为已实现通用依赖求解器。}

\section{基础几何、单位与引用模型}
\label{sec:3-2}

\subsection{基础几何与量纲约定}
\label{subsec:3-2-1}
\outlinehint{定义点、方向、局部坐标系及模型单位；说明长度换算、角度解释和无量纲数据的处理规则。}

\subsection{对象标识、引用角色与几何定位信息}
\label{subsec:3-2-2}
\outlinehint{定义引用对象类别、局部标识、作用角色、几何位置和辅助点等信息，分别承担依赖描述与目标选择定位。}

\subsection{特征依赖与几何快照的组织}
\label{subsec:3-2-3}
\outlinehint{区分统一模型中的构造关系与运行阶段的几何快照和源目标对应记录，交代数据产生时机及后续用途。}

\section{草图与基准面的统一结构}
\label{sec:3-3}

\subsection{草图图元、坐标与构造属性}
\label{subsec:3-3-1}
\outlinehint{列出支承关系、坐标系、图元、局部标识和构造属性，逐项解释用途；区分已有具体图元与样条等未完成范围。}

\subsection{草图约束、尺寸与外部引用}
\label{subsec:3-3-2}
\outlinehint{组织约束类型、参与对象、作用位置、数值与外部来源，说明定义完整不代表全部约束已在三系统恢复。}

\subsection{基准面方法、参考与约束参数}
\label{subsec:3-3-3}
\outlinehint{定义建面方法、参考对象、距离、角度、方向和可选实际几何；明确约束组合才是目标构造的输入。}

\section{拉伸与旋转的统一结构}
\label{sec:3-4}

\subsection{材料操作、终止条件与薄壁共享参数}
\label{subsec:3-4-1}
\outlinehint{定义加料与切除、两端终止类型和参考、偏移启用状态及两侧薄壁偏置；对依赖特征上下文的量纲作明确限定。}

\subsection{拉伸轮廓、方向与作用范围}
\label{subsec:3-4-2}
\outlinehint{给出轮廓引用、实际方向、两端范围、对称总量及附属选项的参数作用表。}

\subsection{旋转轮廓、轴线与角度范围}
\label{subsec:3-4-3}
\outlinehint{分别表达轴的位置和来源、公共正向及两侧角度，保留参考终止所需的方向与辅助信息。}

\section{扫掠的统一结构}
\label{sec:3-5}

\subsection{三类截面及圆管参数表达}
\label{subsec:3-5-1}
\outlinehint{分别定义独立草图引用、内嵌草图和参数圆形截面，说明外半径、内半径及各路径的实际使用范围。}

\subsection{路径、有序引用与起止信息}
\label{subsec:3-5-2}
\outlinehint{定义主路径和相关引用、起点终点、分段顺序及闭合标记，说明缺失数据与明确标记的区别。}

\subsection{截面放置与导引控制参数}
\label{subsec:3-5-3}
\outlinehint{组织截面放置、沿路径的方向控制、导引及薄壁参数；对高级控制保留结构范围与生产范围的差异。}

\section{模型存储与合法性检查}
\label{sec:3-6}

\subsection{序列化与可选参数处理}
\label{subsec:3-6-1}
\outlinehint{说明文本化存储、标识一致性，以及未提供、未启用、零值和未知值的区别；不增加未经实现的格式能力。}

\subsection{参数、引用与组合条件检查}
\label{subsec:3-6-2}
\outlinehint{列出尺寸有效性、必要参考、截面与路径条件等检查，并区分公共合法性与目标系统可执行性。}

\section{本章小结}
\label{sec:3-7}

\outlinehint{归纳本章能够支持的认识，并说明与下一章的衔接；不重复列举所有接口或扩展未证明结论。}

\chapter{五类特征转换方法与关键算法}
\label{chap:4}
\chapterguide{以“问题、输入、步骤、输出、适用条件”组织方法，减少接口名称。关键公式和案例应与真实代码路径及可复现实验对应。}

\section{建模引用的恢复方法}
\label{sec:4-1}

\subsection{源目标对象对应与引用解析}
\label{subsec:4-1-1}
\outlinehint{说明源标识到统一对象、统一对象到目标原生对象的对应，依据历史顺序使用当前已创建的引用。}

\subsection{基于几何信息的目标对象匹配}
\label{subsec:4-1-2}
\outlinehint{利用类型、位置、方向、采样点和必要上下文定位目标边面，处理分段和几何归属变化。}

\subsection{多候选、失效引用与失败处理}
\label{subsec:4-1-3}
\outlinehint{讨论重复或对称几何产生的歧义，明确可用判定、拒绝条件与日志，不把全部候选歧义写成已解决。}

\section{基准面的关联重建方法}
\label{sec:4-2}

\subsection{构造约束与参考组合映射}
\label{subsec:4-2-1}
\outlinehint{按建面方法和参考角色组织目标输入，说明过点、过线、偏置和角度条件的组合与必要性。}

\subsection{偏置方向、角度分支与建后核验}
\label{subsec:4-2-2}
\outlinehint{依据源平面位置和方向选择目标分支，创建后核对实际平面；使用固定辅助点时单独记录关联变化。}

\section{草图几何与约束的恢复方法}
\label{sec:4-3}

\subsection{源目标草图坐标转换}
\label{subsec:4-3-1}
\outlinehint{用源局部坐标到世界坐标、再到目标实际局部坐标的过程推导变换，说明坐标系正交性和单位前提。}

\subsection{圆弧区间与构造几何的恢复}
\label{subsec:4-3-2}
\outlinehint{结合圆心、端点、方向和必要中间点确定同一物理弧段，保留中心线和辅助图元的作用。}

\subsection{约束作用位置与关联恢复}
\label{subsec:4-3-3}
\outlinehint{映射约束参与对象和端点等位置，校正水平、竖直等方向关系；区分驱动尺寸恢复与仅几何复制。}

\section{拉伸方向与终止范围的转换方法}
\label{sec:4-4}

\subsection{实际方向与材料操作的统一}
\label{subsec:4-4-1}
\outlinehint{将源端方向标记转换为明确的有效方向，再映射目标加料和切除路径，避免重复翻转。}

\subsection{单向、双向与对称深度转换}
\label{subsec:4-4-2}
\outlinehint{用两侧材料范围解释深度，对对称总量与单侧量给出换算，并保留双端各自的终止规则。}

\subsection{引用终止与终止偏移处理}
\label{subsec:4-4-3}
\outlinehint{恢复终止对象与偏移，区分有限面和延伸支承面；明确辅助构造、回退定深及拒绝处理对关联的影响。}

\section{旋转方向与角度范围的转换方法}
\label{sec:4-5}

\subsection{旋转轴定位与正转方向判定}
\label{subsec:4-5-1}
\outlinehint{分别恢复轴线位置和公共正向；说明Creo中结合有效轮廓、轴线与草图法向进行判定的依据和退化边界。}

\subsection{双向、对称与正负角度转换}
\label{subsec:4-5-2}
\outlinehint{以实际扫过角域解释两侧参数，使用正90度、反30度对应负30度到正90度的范围作为示例。}

\subsection{参考终止与互补角问题处理}
\label{subsec:4-5-3}
\outlinehint{说明到面、到点的辅助对象与取解要求，明确互补角的未完成范围以及共享偏移量纲的专项检查。}

\section{薄壁材料位置的统一映射方法}
\label{sec:4-6}

\subsection{两侧偏置与材料区间的对应}
\label{subsec:4-6-1}
\outlinehint{用轮廓两侧边界位置表达薄壁，解释单侧、对称及不等量双侧的总厚度和位置差异。}

\subsection{三系统厚度选项映射与限制}
\label{subsec:4-6-2}
\outlinehint{归纳厚度加方向、材料侧与有符号偏置间的映射；明确不等量双侧信息损失等现有边界。}

\section{扫掠截面与路径的联合重建方法}
\label{sec:4-7}

\subsection{路径起点、曲线顺序与方向恢复}
\label{subsec:4-7-1}
\outlinehint{按端点连续性和源选取信息组织路径，说明路径反向、闭合路径与跨对象曲线链的不同处理范围。}

\subsection{截面放置与目标坐标适配}
\label{subsec:4-7-2}
\outlinehint{联合路径起点、切向和截面实际坐标恢复放置，说明截面几何不变仍可能因方向不同生成不同结果。}

\subsection{独立、内嵌与参数圆截面的重建}
\label{subsec:4-7-3}
\outlinehint{分别说明直接引用、补建草图、截面回放以及圆管或双圆截面构造，并检查内半径是否参与生成。}

\subsection{特殊截面补偿与重建策略边界}
\label{subsec:4-7-4}
\outlinehint{解释斜交截面的受限投影补偿及适用条件，区分坐标重表达与真正几何投影；记录重试时合并或方向选项的变化。}

\section{本章小结}
\label{sec:4-8}

\outlinehint{归纳本章能够支持的认识，并说明与下一章的衔接；不重复列举所有接口或扩展未证明结论。}

\chapter{模型转换系统设计与实现}
\label{chap:5}
\chapterguide{交代研究方法怎样进入生产读写流程，突出公共数据与原生执行的边界，避免把各系统列成大量接口清单。}

\section{总体架构与模块分工}
\label{sec:5-1}

\subsection{公共模型层与系统适配层}
\label{subsec:5-1-1}
\outlinehint{说明CADExchange承担统一结构和公共服务，CADProj承担三系统读写、运行调度与展示，给出模块关系图。}

\subsection{原生读取、参数映射与创建执行的分离}
\label{subsec:5-1-2}
\outlinehint{以现有Creo组织说明快照、纯数值映射、写入计划与原生执行的分工，不将三端描述为已经完全相同的结构。}

\section{三系统的读取与重建适配}
\label{sec:5-2}

\subsection{SolidWorks特征读取与选择重建}
\label{subsec:5-2-1}
\outlinehint{概述特征遍历、参数读取、选择对象组织和创建，重点说明参考选择、不同生成分支及附属选项。}

\subsection{Creo元素结构与截面重建}
\label{subsec:5-2-2}
\outlinehint{说明特征参数与截面的读取、独立草图和内嵌截面的组织，以及引用、方向判定与目标创建计划。}

\subsection{NX对象构造与参数适配}
\label{subsec:5-2-3}
\outlinehint{说明对象和构造器路径、草图坐标反投影、对称深度和双向角度适配，以及辅助对象的依赖限制。}

\section{模型读取与顺序重建流程}
\label{sec:5-3}

\subsection{源模型历史、参数与引用提取}
\label{subsec:5-3-1}
\outlinehint{按历史读取五类特征、草图和引用，同时保留单位及必要几何，区分普通图元与构造图元。}

\subsection{目标特征创建与再生状态管理}
\label{subsec:5-3-2}
\outlinehint{按已具备的引用创建目标对象，取得建后实际坐标和可选几何；辅助对象按实际需要插入，不要求节点数完全相同。}

\subsection{运行配置、批处理与模型保存}
\label{subsec:5-3-3}
\outlinehint{交代模板、系统版本、工作单位、执行方式和输出路径；统计完整重建而非仅进程退出或文件生成。}

\section{结果记录、比较与可视化}
\label{sec:5-4}

\subsection{几何快照与源目标对应记录}
\label{subsec:5-4-1}
\outlinehint{说明逐特征几何和对象映射如何产生及复用，保留阶段、来源与状态，支持错误定位。}

\subsection{结构、参数与几何比较的接入}
\label{subsec:5-4-2}
\outlinehint{区分现有几何比较能力与待接回的目标重读、当前字段比较；介绍差异展示，不声称默认完整验证已闭合。}

\subsection{异常、受限处理与日志组织}
\label{subsec:5-4-3}
\outlinehint{列出未支持、读取失败、引用失败、创建失败、语义变化和未完成验证的记录需求，标注哪些仍需完善。}

\section{实现范围与能力边界}
\label{sec:5-5}

\subsection{五类特征的读写覆盖矩阵}
\label{subsec:5-5-1}
\outlinehint{按系统、特征、子类型、读取和创建分别列出范围，保留结构可表达与实际执行的差别。}

\subsection{辅助构造、受限参数与未完成部分}
\label{subsec:5-5-2}
\outlinehint{说明草图驱动尺寸、点与椭圆、特殊终止、不等量薄壁和复杂扫掠等限制，不将局部样例外推为完整支持。}

\section{本章小结}
\label{sec:5-6}

\outlinehint{归纳本章能够支持的认识，并说明与下一章的衔接；不重复列举所有接口或扩展未证明结论。}

\chapter{转换实验设计与结果分析}
\label{chap:6}
\chapterguide{本章为待开展或待统一复现的实验框架。先定义样例、预期、指标与失败口径，再填入实际数据；历史记录不能替代本轮实验结果。}

\section{实验环境、路线与测试样例}
\label{sec:6-1}

\subsection{软件版本、代码版本与容差配置}
\label{subsec:6-1-1}
\outlinehint{固定三套CAD、两仓配套代码、模板、单位、操作系统和比较工具；记录实际配置，不使用未核实的版本号和容差。}

\subsection{同系统基线与跨系统转换路线}
\label{subsec:6-1-2}
\outlinehint{设置三条同系统读取重建基线和六条有向跨系统路线；“往返转换”仅用于实际执行源到目标再到源的试验。}

\subsection{五类特征样例与组合模型}
\label{subsec:6-1-3}
\outlinehint{从现有31组建议案例中筛选可复现样例，安排单因素、参数组合、连续依赖与预期失败模型，先登记预期。}

\section{评价方法与统计指标}
\label{sec:6-2}

\subsection{特征结构、参数与约束检查}
\label{subsec:6-2-1}
\outlinehint{检查特征类型、依赖和参数含义，容许有解释的辅助节点；将尺寸是否驱动、引用是否关联作为独立项目。}

\subsection{曲线、平面与实体的几何比较}
\label{subsec:6-2-2}
\outlinehint{草图检查曲线与位置，基准面检查位置方向，封闭实体检查整体量和双向材料差；处理拓扑分段与接缝差异。}

\subsection{目标原生模型的编辑再生检查}
\label{subsec:6-2-3}
\outlinehint{直接在目标系统修改尺寸、参考面或路径，检查相关特征更新；重新转换源模型不能替代此项。}

\subsection{成功率、失败分类与耗时统计}
\label{subsec:6-2-4}
\outlinehint{分别统计支持范围、有效执行和验证通过；明确分母，区分未支持与失败、无法比较与确定不同，并分阶段记录耗时。}

\section{五类特征的基本转换实验}
\label{sec:6-3}

\subsection{草图与基准面转换实验}
\label{subsec:6-3-1}
\outlinehint{选取倾斜放置、圆弧方向、约束位点、偏置及角度基准面，展示参数、特征树和必要几何比较。}

\subsection{拉伸与旋转转换实验}
\label{subsec:6-3-2}
\outlinehint{覆盖单向、双向、对称、参考终止、薄壁和反向轴等场景，分别核对实际范围、材料位置和引用关系。}

\subsection{扫掠转换实验}
\label{subsec:6-3-3}
\outlinehint{围绕三类截面、非圆轮廓、路径反向和圆管内孔设计样例，区分公共表达范围与各路线实际可执行范围。}

\section{组合模型与编辑能力实验}
\label{sec:6-4}

\subsection{连续依赖模型的转换验证}
\label{subsec:6-4-1}
\outlinehint{使用基准面、草图、成形特征和后续引用构成的连续模型，检查引用恢复及误差传播，列出源目标节点对应。}

\subsection{尺寸与上游引用修改后的再生验证}
\label{subsec:6-4-2}
\outlinehint{修改拉伸深度、旋转角度、基准偏置或扫掠路径，记录初始一致性与修改后的预期结果，定位关联丢失。}

\section{对比实验与综合分析}
\label{sec:6-5}

\subsection{关键处理方法的对照实验}
\label{subsec:6-5-1}
\outlinehint{对比是否进行方向归一化、实际坐标转换、薄壁位置恢复和几何分段处理，控制其他条件；未执行的对照不得写成结果。}

\subsection{转换效率与模型规模影响}
\label{subsec:6-5-2}
\outlinehint{统计读取、统一映射、创建、再生和比较耗时，说明进程启动、首次运行、样本量与模型复杂度。}

\subsection{失败案例、能力边界与适用性分析}
\label{subsec:6-5-3}
\outlinehint{围绕读写不对称、参数简化、引用歧义与目标功能限制解释失败，给出结论适用的路线、子类型和条件。}

\section{本章小结}
\label{sec:6-6}

\outlinehint{归纳本章能够支持的认识，并说明与下一章的衔接；不重复列举所有接口或扩展未证明结论。}

\chapter{总结与展望}
\label{chap:7}
\chapterguide{总结只采用正文和实验能够支持的结论。贡献与不足须对应材料和实际分工，不使用未经比较的“首次”“全面无损”等表述。}

\section{研究工作总结}
\label{sec:7-1}

\subsection{差异分析与统一模型设计总结}
\label{subsec:7-1-1}
\outlinehint{概括参数含义、引用及截面组织的统一表达，指出五类对象之间的共用结构与作用。}

\subsection{转换方法与系统实现总结}
\label{subsec:7-1-2}
\outlinehint{总结坐标和方向处理、引用恢复、材料范围及扫掠联合重建，区分个人工作和协作集成。}

\subsection{实验结论与适用范围总结}
\label{subsec:7-1-3}
\outlinehint{依据第六章真实数据概括可编辑和几何保留程度，明确软件版本、路线、子类型和失败条件；当前留待实验完成。}

\section{研究不足与后续展望}
\label{sec:7-2}

\subsection{草图尺寸、外部关联与终止条件完善}
\label{subsec:7-2-1}
\outlinehint{围绕驱动尺寸、投影关联、复杂参考与不等量薄壁等现有不足提出可执行的改进方向。}

\subsection{复杂扫掠与统一验证体系扩展}
\label{subsec:7-2-2}
\outlinehint{讨论更复杂路径、截面控制和样条的后续研究，以及目标重读、完整参数比较和可复现实验集的完善。}

\chapter*{参考文献}
\markboth{参考文献}{参考文献}
\addcontentsline{toc}{chapter}{参考文献}
\outlinehint{正式文献尚未补齐。围绕几何交换、特征回放、统一模型与转换验证整理实际阅读的论文及官方资料，核对作者、年份、出处和结论。项目代码证据与正式研究文献分开管理。本提纲不添加虚构条目或模板水利学示例文献。}

% 正式文献整理完成后，删除上一行并采用学院确认的参考文献样式。
% 若沿用原项目：从原仓库取得 hhuthesis-numeric.bst 放到本目录，再启用：
% \bibliographystyle{hhuthesis-numeric}
% \bibliography{reference/references}
\appendix
\ctexset{chapter={name={附录,},number={\Alph{chapter}}}}
\chapter{五类特征参数与读写覆盖表}
\outlinehint{汇总参数名称、含义、单位、引用角色，以及各系统的读取和创建范围。每个结论关联固定代码版本和实际样例。}
\chapter{测试样例、转换路线与结果记录表}
\outlinehint{记录样例编号、子类型、源目标系统、预期结果、原生文件、特征对应、几何差异、编辑结果、耗时与失败原因。未执行项保持空白或标为未执行。}
\chapter{主要中间模型结构与关键处理流程}
\outlinehint{给出与正文一致的结构节选、存储示例和必要伪代码。仅列帮助复现实验的关键部分，避免堆叠接口清单或复制全部源码。}

\backmatter
\begin{acknowledgement}
\outlinehint{按实际指导、协作和支持情况撰写。涉及共同完成的功能或数据时，在正文和致谢中分别说明贡献。}
\end{acknowledgement}

\chapter{攻读硕士学位期间取得的成果}
\outlinehint{仅据实填写本人论文、专利、软件著作权、项目及其他成果；核对作者排序、状态和日期。没有的类别不列入，不使用模板示例成果。}

\end{document}
